\subsection{Konfigurationsmodell}
\label{subsec:konfigurationsmodell}

Projektstruktur und Laufzeitkonfiguration sind getrennt. Das Projektmanifest
bestimmt die Features und liegt in \path{project-profile.toml}. Der
Variablenvertrag steht in \path{config/environment.toml}. Für aktive Features gilt die
Priorität CLI-Override, Prozessumgebung, lokale \texttt{.env}, Vertragsdefault.
Auch abgeleitete Werte bleiben nachvollziehbar
\parencite{kleiveist2026configuration}.

Die dritte zentrale Konfigurationsart, \path{config/code-quality.toml}, gehört
nicht zur Laufzeitauflösung. Sie definiert Schema-Version, unterstützte
Endungen, Excludes, Schwellenwerte, Schichtabbildungen und Exceptions für das
Entwicklungs- und CI-Gate \parencite{kleiveist2026quality-config}. Dadurch
bleiben Produktprofil, Laufzeitwerte und Quality-Policy getrennte Verträge.
Die Quality Engine liest das Profilmanifest allerdings nicht selbst; ihre
Komponentenwahl folgt der vorhandenen erzeugten Verzeichnisstruktur.

Abbildung~\ref{fig:configuration-boundaries} zeigt die Sicherheitsgrenzen.
Vite erhält nur die freigegebene \texttt{VITE\_API\_BASE\_URL}; FastAPI
verarbeitet Serverwerte typisiert und maskiert \texttt{DATABASE\_URL}. Das
Tooling reicht nur relevante Werte weiter und entfernt bekannte Secret-Namen
aus der Tauri-Umgebung. Ungültige oder fehlende Werte führen zum kontrollierten
Abbruch, ohne Geheimnisse auszugeben.

\begin{figure}[H]
  \centering
  \resizebox{0.98\textwidth}{!}{%
  \begin{tikzpicture}[node distance=6mm and 22mm]
    \node[casePrimary,minimum width=39mm] (cli) {1. CLI-Override};
    \node[caseBox,minimum width=39mm,below=of cli] (env) {2. Prozessumgebung};
    \node[caseBox,minimum width=39mm,below=of env] (dotenv) {3. lokale \texttt{.env}};
    \node[caseBox,minimum width=39mm,below=of dotenv] (contract)
      {4. Vertragsdefault\\\texttt{environment.toml}};

    \node[caseBox,minimum width=42mm,right=of $(env)!0.5!(dotenv)$] (profile)
      {\texttt{project-profile.toml}\\filtert aktive Variablen};
    \node[casePrimary,minimum width=42mm,below=of profile] (resolver)
      {profilbewusster Resolver\\Priorität und Validierung};

    \node[caseBox,minimum width=39mm,right=of profile] (frontend)
      {öffentlich\\\texttt{VITE\_*}};
    \node[caseBox,minimum width=39mm,below=of frontend] (server)
      {serverseitige Werte};
    \node[caseOptional,minimum width=39mm,below=of server] (secrets)
      {Secrets\\maskiert};
    \node[caseBox,minimum width=39mm,below=of secrets] (errors)
      {ungültig: Abbruch\\ohne Secret-Ausgabe};

    \begin{scope}[on background layer]
      \node[caseGroup,fit=(frontend)(server)(secrets)(errors),
        label={[font=\sffamily\bfseries\scriptsize,text=caseBlue]above:Sicherheitsgrenzen}] {};
    \end{scope}

    \coordinate (inputtop) at ($(cli.east)+(9mm,0)$);
    \coordinate (inputenv) at (inputtop |- env.east);
    \coordinate (inputdotenv) at (inputtop |- dotenv.east);
    \coordinate (inputbottom) at (inputtop |- contract.east);

    \coordinate (outputtop) at ($(frontend.west)+(-9mm,0)$);
    \coordinate (outputserver) at (outputtop |- server.west);
    \coordinate (outputsecrets) at (outputtop |- secrets.west);
    \coordinate (outputbottom) at (outputtop |- errors.west);

    \begin{scope}[on background layer]
      \draw[draw=caseLine,line width=0.8pt] (cli.east) -- (inputtop);
      \draw[draw=caseLine,line width=0.8pt] (env.east) -- (inputenv);
      \draw[draw=caseLine,line width=0.8pt] (dotenv.east) -- (inputdotenv);
      \draw[draw=caseLine,line width=0.8pt] (contract.east) -- (inputbottom);
      \draw[draw=caseLine,line width=0.8pt] (inputtop) -- (inputbottom);
      \draw[caseFlow] (inputtop |- resolver.west) -- (resolver.west);
      \draw[caseFlow] (profile) -- (resolver);

      \draw[draw=caseBlue,line width=0.9pt] (resolver.east) -- (outputserver);
      \draw[draw=caseBlue,line width=0.9pt] (outputtop) -- (outputsecrets);
      \draw[caseFlow] (outputtop) -- (frontend.west);
      \draw[caseFlow] (outputserver) -- (server.west);
      \draw[caseFlow] (outputsecrets) -- (secrets.west);
      \draw[caseSecondaryFlow] (outputsecrets) -- (outputbottom) -- (errors.west);
    \end{scope}
  \end{tikzpicture}}
  \caption{Konfigurationspriorität und Sicherheitsgrenzen}
  \label{fig:configuration-boundaries}
  \smallskip
  {\footnotesize Quelle: Eigene Darstellung auf Grundlage von
  \textcite{kleiveist2026configuration}.\par}
\end{figure}

\beginnerinput{workspace/statement/300_architektur/360}
